% Direct-authority diplomatic French for printed p.201.
\clearpage
\SourcePage{201}

\LayerTarget{tohoku:proof:theorem-5.2.1:p1}
\noindent atement du fait que tout $G$-faisceau abélien peut se plonger dans
un $G$-faisceau abélien flasque (comme nous avons vu au numéro précédent),
qui annule donc les $H^q(X,A)$ pour $q>0$.

\LayerTarget{tohoku:construction:edge-homomorphisms-5.2}
Ces suites spectrales donnent lieu tout d'abord à des
``edge-homomorphismes'' importants :

\LayerTarget{tohoku:equation:5.2.6}
\[
H^n(Y,A^G)\to H^n(X;G,A)
\to H^0(Y,\mathbf H^n(G,A)),
\tag{5.2.6}
\]

\LayerTarget{tohoku:equation:5.2.7}
\[
H^n(G,H^0(X,A))\to H^n(X;G,A)
\to H^n(X,A)^G.
\tag{5.2.7}
\]

\LayerTarget{tohoku:construction:edge-map-interpretation}
La deuxième application (5.2.6) s'interprète immédiatement si on note qu'en
vertu de lemme 3.7.2, $\mathbf H^n(G,A)$ est le faisceau sur $Y$ associé au
préfaisceau formé par les $H^n(f^{-1}(U);G,A)$. Le composé du premier
homomorphisme de la première ligne avec le deuxième homomorphisme de la
deuxième ligne est l'homomorphisme $f^*$ associé à l'injection naturelle
(5.1.2) de $f^{-1}(A^G)$ dans $A$ (cf. 3.2).

\LayerTarget{tohoku:result:five-term-exact-sequences-5.2}
Les suites spectrales du théorème 5.2.1 définissent aussi deux suites exactes
à 5 termes :

\LayerTarget{tohoku:equation:5.2.8}
\[
\tag{5.2.8}
\begin{aligned}
0\to H^1(Y,A^G)\to H^1(X;G,A)
&\to H^0(Y,\mathbf H^1(G,A))\to H^2(Y,A^G)\\
&\to H^2(X;G,A),
\end{aligned}
\]

\LayerTarget{tohoku:equation:5.2.9}
\[
\tag{5.2.9}
\begin{aligned}
0\to H^1(G,\Gamma A)\to H^1(X;G,A)
&\to H^1(X,A)^G\to H^2(G,\Gamma A)\\
&\to H^2(X;G,A).
\end{aligned}
\]

\LayerTarget{tohoku:construction:third-spectral-sequence}
Il y a une troisième façon de représenter $\Gamma_X^G$ comme un foncteur
composé, savoir $\Gamma_X^G=\Gamma_Y^Gf_*$, et le théorème 2.4.1 est encore
appliquable en vertu de proposition 5.1.3. Donc, le foncteur gradué
$(H^n(X;G,A))$ est aussi l'aboutissement d'une troisième suite spectrale dont
le terme initial est

\LayerTarget{tohoku:formula:third-spectral-initial-term}
\[
E_2^{p,q}=H^p(Y;G,R^qf_*(A)).
\]
\LayerTarget{tohoku:equation:5.2.10}
Cette suite spectrale définit des homomorphismes fonctoriels
\[
H^n(Y;G,f_*(A))\to H^n(X;G,A).
\tag{5.2.10}
\]
\LayerTarget{tohoku:remark:third-spectral-sequence-pathology}
L'intérêt de cette suite spectrale est limitée, car il apparait qu'elle est
pathologique dans les cas où elle n'est pas triviale, c'est-à-dire quand on
n'est pas dans les conditions de la proposition suivante :

\LayerTarget{tohoku:proposition:5.2.2}
\textsc{Proposition 5.2.2.} \textit{Supposons le $G$-faisceau
abélien $A$ tel que $R^qf_*(A)=0$ pour $q>0$. Alors les homomorphismes
\textup{(5.2.10)} sont des isomorphismes, de plus les deux suites spectrales du
théorème \textup{5.2.1} s'identifient alors aux deux suites spectrales correspondantes
pour le $G$-faisceau $f_*(A)$ sur $Y$, et on a}

\LayerTarget{tohoku:equation:5.2.11}
\[
\mathbf H^q(G,A)=\mathbf H^q(G,f_*(A)).
\tag{5.2.11}
\]

\LayerTarget{tohoku:proof:proposition-5.2.2:p1}
La première assertion résulte aussitôt de la forme du terme initial de la
suite spectrale $E$ définissant (5.2.10). Les autres assertions résultent de
la définition explicite des fermes qu'il s'agit de comparer, si on note que
pour toute résolution injective $\mathrm C(A)$ de $A$ dans ${\mathbf C}^{X(G)}$, le complexe
$f_*(C(A))$ est une \textit{résolution} de $f_*(A)$ (c'est ce que signifie
l'hypothèse $R^qf_*(A)=0$ pour $q>0$!) qui est injective en vertu de
proposition 5.1.3. --- Notons que nous avions implicitement défini les
$R^qf^*(A)$ comme dérivés droits de $f_*$ considéré comme

\TohokuContinuation{202}
